\section{Discussion and Limitations}\label{sec:disc}

\noindent\textbf{What the Factoring Buys.} Because the MCL core is written against one \code{MeasurementModel} interface, the 2D and 3D backends share the motion model, weighting, KLD resampling, estimate, and every test of those components; adding a backend means adding one likelihood function. Because the cores are middleware-free, every equation of \Cref{sec:method} is exercised by deterministic tests and by the closed-loop drivers of \Cref{sec:exp} without a ROS installation, which is also what makes the evaluation reproducible bit for bit from the repository.

\noindent\textbf{Cost Analysis.} The per-update cost of the MCL backends is linear in the particle count and in the number of scored beams or points (\Cref{eq:laser,eq:ndt}); KLD adaptation keeps the former at the 500-particle floor once the posterior is concentrated (\Cref{fig:mcl}). The branch-and-bound matcher's cost is data-dependent through pruning: on the synthetic room a full-map query took a median of 37.65\,ms (min 8.77\,ms, 95th percentile 106.07\,ms) on the machine that produced the committed data, and it runs only when relocalization is requested, not every scan. The ESKF step is a fixed $15\times15$ covariance propagation per IMU sample.

\noindent\textbf{Limitations.} \sname{} is an architecture paper, and its scope is deliberately bounded.
(i)~\emph{Planar state in the MCL backends.} Both particle-filter backends estimate $(x,y,\theta)$; \code{ndt3d} evaluates a full 3D Gaussian-voxel map from that planar pose at a fixed base height, which suits wheeled robots on locally flat floors but does not track roll, pitch, or vertical motion. Full $SE(3)$ state exists only in \code{fusion3d}.
(ii)~\emph{2D-only global relocalization.} \Cref{alg:bbs} operates on a 2D occupancy pyramid; \code{ndt3d} and \code{fusion3d} rely on an operator-supplied or GNSS-derived initial pose and local tracking.
(iii)~\emph{No time-offset or extrinsic estimation.} The ESKF treats the LiDAR/GNSS-to-IMU time offset and the LiDAR--IMU extrinsic as known inputs.
(iv)~\emph{Loosely-coupled LiDAR update.} In \code{fusion3d}, PCL NDT produces a 6-DoF pose measurement that the ESKF fuses as a loosely-coupled update, which keeps the filter core independent of the registration engine but forgoes the accuracy of tightly-coupled formulations such as FAST-LIO \citep{xu2021fastlio}; the implementation also omits the reset Jacobian of \citet{sola2017quaternion}.
(v)~\emph{Single static map.} All backends assume one prior map of a static world, loaded once, with no online map update or change detection.
(vi)~\emph{No field-accuracy evaluation.} The evaluation of \Cref{sec:exp} is deterministic unit and integration testing plus a seeded synthetic study whose worlds and sensor-noise models match the filters' assumptions by construction, and it runs no external baseline. It establishes that each core behaves as specified in closed loop, not how accurate \sname{} is on a physical robot. We have not measured trajectory error on physical or public datasets; the planned evaluation will report absolute and relative trajectory error (ATE/RPE) following the TUM RGB-D conventions \citep{sturm2012benchmark} on the per-backend \code{map}$\to$\code{odom} estimates. Until then we make no accuracy claim.

\noindent\textbf{Future Work.} Extending the MCL core to a 6-DoF particle state or coupling it to an external attitude source, a 3D branch-and-bound or place-recognition front end for the 3D backends, joint time-offset and extrinsic estimation in the ESKF, and the public-dataset evaluation above are the next steps on the roadmap.
